\section{Initial Value Problems (IVPs)}

\newcommand{\by}{\mathbf{y}}
\newcommand{\bp}{\mathbf{p}}
\begin{frame}
$$\begin{cases}
    \dot{\by} = f(t, \by; \bp),  \\
    \by(t_0) = \by_0,
\end{cases}$$
$\by(t)$: state vector, $\bp$: vector of fixed parameters. Goal: given \underline{many} different $\by_0, \bp$, compute $\by(t_f)$ for each system.

\begin{itemize}
    \item Important application: astrochemistry. We need to solve a large IVP for each spatial grid cell, for each (outer loop) time step. This is a major computational bottleneck in these simulations; it is necessary to have surrogate/fast batched solver.
\end{itemize}
    
\end{frame}

\begin{frame}{Time Stepping}
The simplest time-stepping algorithm is \textbf{forward Euler}:
\begin{gather}
\by_{k+1} = \by_k + \Delta t \cdot f(t_k, \by_k; \bp) \tag{FE} \\
\by_{k+1} = \by_k + \Delta t \cdot f(t_k, \by_{k+1}; \bp)\label{eq:BE} \tag{BE} 
\end{gather}


\begin{itemize}
    \item This is a \textbf{first-order method}: as $\Delta t \to 0$, the error of the solution decreases at a linear rate.
    \item Extremely simple to implement, assuming the time step is fixed.
    \item Extremely easy to implement and parallelize.
    
\end{itemize}
\end{frame}

\begin{frame}{Adaptive Time-Stepping}
In most applications, we would like to change the time step to match the dynamics. This is done by considering two time-steppers of different order and comparing the difference. For example: \textbf{RK23} (Bogacki, Shampine 1997):
\small{
\begin{align*}
    k_1 &= f(t_k, \by_k) \\
    k_2 &= f\left(t_k + \frac{\Delta t}{2}, \by_k + \frac{\Delta t}{2}k_1\right) \\
    k_3 &= f\left(t_k + \frac{3\Delta t}{4}, \by_k + \frac{3\Delta t}{4}k_2\right) \\
    \by_{k+1}^{(3)} &= \by_k + \frac{\Delta t}{9}(2k_1 + 3k_2 + 4k_3) \\
    t_{k+1} &= t_k + \Delta t \\
    k_4 &= f\left(t_{k+1}, y_{k+1}^{(3)}\right) \\
    e_{k+1} &= \frac{\Delta t}{72}(-5k_1 + 6k_2 + 8k_3 - 9k_4).
\end{align*}}
\end{frame}
\begin{frame}
\begin{itemize}
    \item If $\|e_{k+1}\| > \epsilon_{\mathrm{tol}}$, we shrink the time step and recompute. In either case we compute a new time step to use for the next:
\end{itemize}
$$\Delta t^{(k+1)} = \Delta t^{(k)} * \left(\frac{\epsilon_{\mathrm{tol}}}{\|e_{k+1}\|}\right)^{0.2}.$$
\textbf{This is more difficult to parallelize on a GPU since some threads may desynchronize.}
\begin{itemize}
    \item If each ODE is solved on a single thread, some ODEs may reach $t_f$ faster than others. Hence many threads will remain idle while the others finish.
\end{itemize}
\end{frame}
